\section{Succinct learning by order medians}\label{sec:fast}
The real-grid structure permits a learner that does not construct the truth table, its cap mixture, or a radial-isotropic transform. Its negative labeled measure is supported on one unknown affine line. Integer order, rather than pairwise sampling, identifies that line unless one negative point already carries enough mass for a star reduction. This resolves the all-light case with a quadratic, rather than linear, logarithmic query bound.

Write $\mu^-(E)=\Pr_D[p\in E,\ h^*(p)=-1]$, $\beta=\mu^-(X)$, and $\zeta=\epsilon/4096$. We analyze answers of error at most $\zeta$. In a finite-bit implementation request raw SQ tolerance $\zeta/2$ and round each response with additional error at most $\zeta/2$, using $O(\log(1/\epsilon))$ bits after the binary point. This makes the algorithm a finite-bit computation without assuming exact real arithmetic. Points have coordinates $(x,y)$; the label in an SQ is a separate variable.

\begin{samepage}
\begin{lemma}[Noisy integer medians]\label{lem:median}
Suppose $\beta>0$, $|\widehat\beta-\beta|\le\zeta$, and $\beta\ge512\zeta$. Let $f$ be integer-valued in $[L,H]$ on the negative support. Binary search using the query $\mu^-\{f\le t\}$ and threshold $\widehat\beta/2$ returns $v\in[L,H]$ with
\[
 \mu^-\{f\le v\}\ge\beta/2-3\zeta/2,
 \qquad
 \mu^-\{f<v\}<\beta/2+3\zeta/2.
\]
If the answer $\widehat m$ to $m=\mu^-\{f=v\}$ is at most $\widehat\beta/8$, then
\begin{equation}\label{eq:balancedmedian}
 \beta/2-3\zeta/2\le\mu^-\{f\le v\}
 <5\beta/8+21\zeta/8.
\end{equation}
Otherwise $m>\beta/9$.
\end{lemma}
\end{samepage}
\begin{proof}
Start with the bracket $[L-1,H]$, regarding the lower endpoint as a negative comparison and the upper as a positive comparison. The corresponding true masses are zero and $\beta$, so both satisfy the claimed endpoint inequalities. At an interior query, an answer at least $\widehat\beta/2$ implies true mass at least $\beta/2-3\zeta/2$; an answer below it implies true mass below $\beta/2+3\zeta/2$. Preserve the respective endpoint inequalities during binary search. They remain valid even if different valid answers do not form a monotone sequence. At termination the endpoints are $v-1,v$, proving the first two assertions.

In the light-atom case, $m\le\widehat\beta/8+\zeta\le\beta/8+9\zeta/8$. Add this to the upper bound on $\mu^-\{f<v\}$ to obtain \eqref{eq:balancedmedian}. In the other case $m>\widehat\beta/8-\zeta\ge\beta/8-9\zeta/8>\beta/9$, since $\beta\ge512\zeta$.
\end{proof}

For a point $p=(p_x,p_y)$, its template star contains the slopes
\[
 1\le a\le t(p):=\min\{N,\lfloor(p_y-1)/p_x\rfloor\},
 \qquad b=p_y-ap_x.
\]
The upper intercept bound is automatic. Distinct star lines intersect only at $p$. Say that a flip matrix is \emph{star-regular} if for every point with $t(p)\ge64\lceil\log_2(2N^3)\rceil$, at least $t(p)/4$ of these rows are negative at $p$.

\begin{samepage}
\begin{theorem}[Succinct median learner]\label{thm:fast}
Suppose entries of a real-grid incidence flip $A$ can be evaluated in time $E(n)$, where $n=\log_2(2N^3)$. There is a deterministic improper distribution-independent SQ learner of error at most $\epsilon$, using $O(n^2)$ queries at raw tolerance $\epsilon/8192$, with total computation polynomial in $(n,1/\epsilon,E(n))$ and an efficiently evaluable output. If $A$ is star-regular, there is a proper randomized learner with
\[
 m=O(n^2+1/\epsilon+\log(1/\epsilon)),\qquad
 \tau=\epsilon/8192,
\]
with the same polynomial-time property and expected error at most $\epsilon$ against every valid oracle.
\end{theorem}
\end{samepage}
\begin{proof}
Query $\beta$, obtaining $\widehat\beta$. If $\widehat\beta\le\epsilon/4$, return the all-positive row $(a,b)=(1,2N^2)$. Its loss is at most $\epsilon/4+\zeta$. Otherwise
\[
 \beta>\epsilon/4-\zeta>\epsilon/8,
 \qquad \beta>1023\zeta>512\zeta.
\]
All subsequent negative-measure queries are unconditioned indicators and have error at most $\zeta$.

First take a noisy median $v_x$ of $x$. If its atom is heavy in Lemma~\ref{lem:median}, it consists of one negative point: a nonvertical target line has at most one point with a given $x$ coordinate. Binary queries $\mu^-\{x=v_x,y\le t\}$ locate its $y$ coordinate. Each such query has true mass either zero or more than $\beta/9$. Threshold $\widehat\beta/32$ distinguishes these cases despite error $\zeta$: it exceeds $\zeta$ and is less than $\beta/9-\zeta$. We then use the heavy-point completion described below.

Otherwise fix the prefix $P=\{x\le v_x\}$, whose mass satisfies \eqref{eq:balancedmedian}. Maintain an integer interval containing the unknown slope $a^*\in[1,N]$. At its midpoint $a$, take a noisy median $v$ of
\[
 f_a(x,y)=y-ax,
 \qquad 1-N^2\le f_a(x,y)\le2N^2.
\]
On the target line this is $(a^*-a)x+b^*$. Query the mass $m$ of its median atom.

Suppose that atom is heavy. Query its intersection mass $l$ with $P$. If both $\widehat l\ge\widehat\beta/8$ and $\widehat m-\widehat l\ge\widehat\beta/8$, then $a=a^*$ and $v=b^*$. Indeed, for $a\ne a^*$, $f_a$ is injective on the target line, so its atom is a single point and lies wholly on one side of $P$. One of the true split masses is zero. Its directly observed value is at most $\zeta$, or its difference of two observed values at most $2\zeta$, both less than $\widehat\beta/8$. Conversely, when $a=a^*$, the atom has mass $\beta$ and the median value is exactly $b^*$. The two true split masses are at least $\beta/2-3\zeta/2$ and $3\beta/8-21\zeta/8$ by \eqref{eq:balancedmedian}. Their observed lower bounds exceed $\widehat\beta/8$, using $\beta>512\zeta$. Thus the test accepts precisely the constant-residual case. Return the exact target row in that case. If the split test fails, the heavy atom is a single negative point; locate its $x$ coordinate by binary queries restricted to $\{f_a=v\}$, then put $y=ax+v$ and use the heavy-point completion.

It remains to handle a light residual atom. The set $Q=\{f_a\le v\}$ also satisfies \eqref{eq:balancedmedian}. If $a^*>a$, the restrictions of $P,Q$ to the target line are prefixes of the same order. Consequently
\[
 \mu^-(P\cap Q)\ge\beta/2-3\zeta/2.
\]
If $a^*<a$, they are a prefix and a suffix. Their union covers the entire negative support whenever they overlap. Thus, whether or not they overlap,
\[
 \mu^-(P\cap Q)
 \le\max\{0,\mu^-(P)+\mu^-(Q)-\beta\}
 \le\beta/4+21\zeta/4.
\]
Query this intersection and compare its answer to $3\widehat\beta/8$. In the first case the answer is at least $\beta/2-5\zeta/2>3\widehat\beta/8$; in the second it is at most $\beta/4+25\zeta/4<3\widehat\beta/8$. Both inequalities follow from $\beta>512\zeta$. Equality of slopes is impossible in the light case, because the entire negative mass would be one residual atom. Retain the appropriate strict half of the slope interval. After $O(\log N)$ stages only $a^*$ remains. A median query search for its constant residual gives $b^*$, so return $h_{a^*,b^*}$. Each slope stage uses $O(\log N)$ queries. This is the all-light algorithm.

\paragraph{Heavy-point completion.}
We have located an actual negative point $p$. Query the negative tail
\[
 \beta_p=\mu^-(X\setminus\{p\}).
\]
If its answer exceeds $\epsilon/4$, then $\beta_p>\epsilon/4-\zeta$. For a half-interval of slopes in the star at $p$, query negative mass away from $p$ on the union of those lines. Every negative point other than $p$ determines the same unique target slope, so this query is exactly zero or $\beta_p$. Threshold $\epsilon/8$ distinguishes the cases. Binary search identifies the target in $O(\log N)$ queries. The union is succinctly evaluable: for a query point $q\ne p$, test $q_x\ne p_x$, divisibility of $q_y-p_y$ by $q_x-p_x$, and membership of the quotient in the chosen slope interval. This requires no enumeration and no evaluation of flips.

If the tail answer is at most $\epsilon/4$, the improper learner returns the point predictor that is negative only at $p$. Its loss is $\beta_p\le\epsilon/4+\zeta$. This proves the deterministic improper statement for every evaluable flip matrix.

For proper output, put
\[
 C_0=\max\{64\lceil\log_2(2N^3)\rceil,\lceil128/\epsilon\rceil\}.
\]
If $t(p)\le C_0$, enumerate its slopes, discard rows positive at $p$ by an entry evaluation, and query the full loss of each remaining row until an answer is at most $\epsilon/2$. The target is among them and must pass. A passing row has actual loss at most $\epsilon/2+\zeta$.

If $t(p)>C_0$, star-regularity supplies at least $t(p)/4$ rows negative at $p$. The sets of negative points of different star rows outside $p$ are disjoint. At most $8/\epsilon$ rows can therefore have $D$-mass greater than $\epsilon/8$ on their outside negative sets. At least $t(p)/4-8/\epsilon\ge t(p)/8$ rows are negative at $p$ and have outside negative mass at most $\epsilon/8$. Every such row has loss at most
\[
 \beta_p+\epsilon/8\le3\epsilon/8+\zeta,
\]
and its loss answer is at most $3\epsilon/8+2\zeta<\epsilon/2$, regardless of the oracle. Sample
\[
 J=\lceil16\ln(4/\epsilon)\rceil
\]
independent star slopes as follows: draw $q=\lceil\log_2t(p)\rceil$ unbiased bits, interpret them as an integer, reduce modulo $t(p)$, and add one. Every slope has probability at least $1/(2t(p))$, so a guaranteed-passing row has probability at least $1/16$ in each trial. Evaluate its label at $p$ and test negative rows as above. The chance that no guaranteed-passing row is sampled is at most $(1-1/16)^J\le\epsilon/4$. On that event return the all-positive row. All other outputs have error at most $\epsilon/2+\zeta$. The expected loss is at most $\epsilon/2+\zeta+\epsilon/4<\epsilon$.

There are $O(n^2)$ median queries, $O(n)$ locating and star-search queries, and at most $O(n+1/\epsilon+\log(1/\epsilon))$ candidate tests. Every predicate uses polynomially many $n$-bit integer operations and at most one entry evaluation. The output is either a row index with its evaluation algorithm or, in the improper case, a point predictor. The finite answer rounding specified before Lemma~\ref{lem:median} completes the bit-complexity assertion.
\end{proof}

\begin{samepage}
\begin{lemma}[Regular stars for almost every key]\label{lem:regularstars}
For independent fair incidence flips, the probability of failure of star-regularity is at most $e^{-3n}$, where $n=\log_2(2N^3)$. Under the PRF assumption and parameters of Theorem~\ref{thm:prf}, failure occurs for only a negligible fraction of keys.
\end{lemma}
\end{samepage}
\begin{proof}
At a fixed point with $t$ incident rows, the number of negatives is a sum of $t$ independent fair bits. For $\lambda=\ln3$,
\[
 \Pr[Z<t/4]\le e^{\lambda t/4}\left(\frac{1+e^{-\lambda}}2\right)^t
 =\left(\frac{2}{3^{3/4}}\right)^t\le e^{-t/16}.
\]
For the last elementary inequality, $\ln(3^{3/4}/2)>1/16$; for example $3^3/2^4=27/16>e^{1/4}$, since $e^{1/4}\le(1-1/4)^{-1}=4/3<27/16$. If $t\ge64\lceil n\rceil$, the failure probability is at most $e^{-4n}$. There are at most $2^n$ points, so a union bound gives $e^{-(4-\ln2)n}\le e^{-3n}$.

A distinguisher can read every incidence bit and test all stars in time $2^{O(n)}$. This is below $2^{\sigma^\delta}$ for the stated key length. The random-function bound plus PRF indistinguishability proves the claim. Polynomial growth of $\sigma$ in $n$ makes $e^{-3n}$ negligible in the key length as well.
\end{proof}

\begin{samepage}
\begin{corollary}[Conditional polynomial-time, seed-known separation]\label{cor:fastprf}
Under the PRF assumption, for each fixed $c\ge1$ all but a negligible fraction of keys in Theorem~\ref{thm:prf} define classes with
\[
 \dc(H_{A_s})\ge\frac{n^c}{30\log_2n},\qquad
 m=O_\epsilon(n^2),\qquad\tau=\Omega_\epsilon(1),
\]
and a proper learner running in $\operatorname{poly}_{c,\epsilon}(n)$ time given the key. Every key has the corresponding deterministic improper polynomial-time learner, independently of star-regularity.
\end{corollary}
\end{samepage}
\begin{proof}
Intersect the high-rank key event from Theorem~\ref{thm:prf} with the star-regular event from Lemma~\ref{lem:regularstars}. Their two negligible exceptional probabilities have negligible sum. Apply Theorem~\ref{thm:fast} and the polynomial-time entry evaluator.
\end{proof}

Taking $c>2$ refutes the linear query/tolerance implication even when the learner must run in time polynomial in $n$. Taking $c$ larger than twice a proposed polynomial degree refutes each fixed-degree polynomial in $(m,1/\tau,n)$ within this computational model. The key length and runtime exponents grow with $c$, so this does not refute a bound whose right side also charges an arbitrary polynomial of the actual runtime or expanded description size. The $O(n)$-query, polynomial-in-$n$ implementation requested by the stronger stretch target remains open; the proved implementation uses $O(n^2)$ queries.
